Lors d'une fête communale, $n$ personnes sont candidates pour participer à
l'organisation.
\begin{enumerate}
  \item Combien y a-t-il de façons de choisir $4$ personnes et de leur
  attribuer les rôles suivants : président, vice-président, secrétaire,
  trésorier ?
  \item Combien y a-t-il de façons de choisir simplement un groupe de $4$
  personnes, sans rôle ?
  \item Retrouver à partir de la question 1 la formule
  \[
  \binom{n}{4} = \frac{n!}{4!(n-4)!}.
  \]
  \item Calculer $\binom{11}{4}$.
  \item Même question si l'on choisit $4$ personnes parmi $11$, mais que deux
  rôles seulement sont distincts : un président et un secrétaire, les deux
  autres membres n'ayant pas de rôle particulier.
\end{enumerate}
